\documentclass[11pt,a4paper]{article}
\usepackage[margin=1.05in]{geometry}
\usepackage{amsmath,amssymb,amsthm}
\usepackage{microtype}
\usepackage[hidelinks]{hyperref}
\hypersetup{pdftitle={A midpoint identity for the complementary torus order at discriminant -3},pdfauthor={Jack Pickett}}

\newtheorem{lemma}{Lemma}
\newtheorem{remark}{Remark}

\title{A midpoint identity for the complementary torus order at discriminant \(-3\)}
\author{Jack Pickett \textperiodcentered\ 27th September 2026}
\date{}

\begin{document}
\maketitle
\thispagestyle{empty}

\begin{abstract}
For an odd prime \(p\neq 3\), write \(\chi_3(p)=(-3/p)\) and \(m_0(p)=p+\chi_3(p)\). Then
\[
\frac1{m_0(p)}-\frac12\Bigl(\frac1{p-1}+\frac1{p+1}\Bigr)=-\frac{\chi_3(p)}{p^2-1}.
\]
The identity is elementary. Summing the right-hand side gives an explicit convergent series for the offset between the selected neighbour and the midpoint of \(\{p-1,p+1\}\). The same character is the Kronecker symbol that distinguishes the two torus orders of Williams' \(p+1\) test at \(D=-3\). No asymptotic for the associated hire-count is claimed, nor a relation to \(L(s,\chi_3)\) beyond the identification \(\chi_3=(-3/\cdot)\).
\end{abstract}

\section{The door}

Let \(p>3\) be prime and let \(\chi_3\) be the non-principal Dirichlet character modulo \(3\), equivalently the Kronecker symbol \(\bigl(\frac{-3}{\cdot}\bigr)\). Set
\[
m_0(p)=p+\chi_3(p)=
\begin{cases}
p+1 & \text{if }p\equiv 1\pmod{3},\\
p-1 & \text{if }p\equiv 2\pmod{3}.
\end{cases}
\]
Thus \(m_0(p)\) is one of the two even neighbours of \(p\). Williams' \(p+1\) method at discriminant \(D=-3\) uses the complementary order
\[
m_1(p)=p-\chi_3(p).
\]
The pair \(\{m_0,m_1\}\) is \(\{p-1,p+1\}\). The hire graph keeps \(m_0\) and discards \(m_1\). Same quadratic field as the Eisenstein integers and as the CM curve \(y^2=x^3+B\); different group.

\section{The identity}

\begin{lemma}
\label{lem:midgap}
For every prime \(p>3\),
\[
\frac1{m_0(p)}-\frac12\Bigl(\frac1{p-1}+\frac1{p+1}\Bigr)=-\frac{\chi_3(p)}{p^2-1}.
\]
\end{lemma}

\begin{proof}
The midpoint reciprocal is \(p/(p^2-1)\). If \(p\equiv 1\pmod{3}\), then \(\chi_3(p)=+1\) and \(m_0(p)=p+1\), so the left-hand side is
\[
\frac1{p+1}-\frac{p}{p^2-1}=\frac{p-1-p}{(p+1)(p-1)}=-\frac1{p^2-1}.
\]
If \(p\equiv 2\pmod{3}\), then \(\chi_3(p)=-1\) and \(m_0(p)=p-1\), so the left-hand side is
\[
\frac1{p-1}-\frac{p}{p^2-1}=\frac{p+1-p}{(p-1)(p+1)}=+\frac1{p^2-1}.
\]
In both cases this equals \(-\chi_3(p)/(p^2-1)\).
\end{proof}

\begin{remark}
Equivalently, the midpoint is \(p/(p^2-1)\) and the gap is the unit \(\pm 1/(p^2-1)\). Adding them recovers \(1/m_0(p)\).
\end{remark}

\section{The offset}

Summing Lemma~\ref{lem:midgap} over primes \(p>3\) up to \(X\) gives
\[
\sum_{5\le p\le X}\frac1{m_0(p)}-\frac12\sum_{5\le p\le X}\Bigl(\frac1{p-1}+\frac1{p+1}\Bigr)=-\sum_{5\le p\le X}\frac{\chi_3(p)}{p^2-1}.
\]
The right-hand side converges absolutely as \(X\to\infty\). Its value is
\[
-\sum_{p>3}\frac{\chi_3(p)}{p^2-1}\approx 0.02598841679.
\]
There is no reason to expect a form in \(\pi\) and \(\sqrt{3}\): \(L(1,\chi_3)=\pi/(3\sqrt{3})\) and \(L(3,\chi_3)=4\pi^3/(81\sqrt{3})\) are odd arguments of an odd character; the present series is a character-prime-zeta at even powers. The series is left as the series.

The unsigned neighbour sums \(\sum 1/(p-1)\) and \(\sum 1/(p+1)\) each diverge like \(\tfrac12\log\log X\). The identity says the selected door tracks the midpoint of those two diverging curves at a frozen gap. The signed series \(\sum\chi_3(p)/m_0(p)\) is a different object (cancellation at \(1/p\), numerical value \(\approx-0.241\)) and is not used here.

\section{Refusal}

Lemma~\ref{lem:midgap} does not yield an asymptotic for
\[
H(X)=\#\{q\text{ prime}:5\le q\le X,\ \exists\,p\le X,\ q\mid m_0(p)\}.
\]
It is not an elliptic-curve method. It records one elementary relation between the complementary torus order at \(D=-3\) and the two neighbours of \(p\).

\end{document}
